\documentclass{handout}

\assignment{Homework 4}
\title{Forces and Newton's laws}

\begin{document}
\maketitle
\practiceinstructions

\section{Inertia and free body diagrams}

\begin{questions}
\question \textbf{First law}: In each of the following cases, explain whether there is a net force on the object.
\begin{parts}
    \part A bicycle hub lies stationary, connected to 28 spokes under tension.
    \part An elevator departing from or arriving at a floor.
    \part An elevator moving at constant speed between floors.
    \part A person in a roller coaster going over the crest of a hill at constant speed.
    \part A raindrop falling at constant velocity (its terminal velocity).
\end{parts}
\question \textbf{Kitesurfer}: A kitesurfer on lake Michigan of mass 70 kg experiences a gust of wind which takes him airborne. At the moment he leaves the water, he experiences three forces: his own weight $\vec{W}$, a 1000 N tension force $\vec{T}$ 53$^\circ$ above the horizonal from the kite string, and a 200 N drag from the air $\vec{D}$ pushing backward (to the left).  
\begin{parts}
    \part Draw a free body diagram of the kitesurfer with all three forces labelled.
    \part Find the net force on the kitesurfer in component form.
    \part What is the magnitude of the kitesurfer's acceleration? 
    \part What angle above the horizontal is his acceleration?
\end{parts}

\end{questions}

\section{Friction}


\begin{questions}
\question \textbf{Friction}: You are pulling a block (mass $m$) with force $F$ (magnitude) to the right of the page via a rope towed over your shoulder. As a result, the rope forms an angle of $\theta$ with the horizontal. It slides with coefficient of friction $\mu_k$.
\begin{parts}
    \part Draw a free body diagram of the forces on the block.
    \part Find the normal force which the ground exerts to keep the block from falling through the ground.
    \part Find an expression for the net force on the block. \part Explain why it is not possible for this expression to return a negative value (force to the left).
    \part You are pulling the block with a constant velocity, applying a force of 200 N at an angle of 30$^\circ$. The block's mass is 30 kg. What is the coefficient of friction $\mu$?
    \part Explain why your answer to part b is better than how you would have done the problem if you had all the numbers from part d up front.
\end{parts}

\question \textbf{Critical angle}: If something is sitting on an inclined surface and held up only by friction, there is \emph{always} a maximally inclined angle at which the object will slide off. 
\begin{parts}
\part Normal force can support \emph{as much force as needed} (until you're destroying the surface) to prevent an object from accelerating in the direction of a surface with which it is in contact. Find an expression for the normal force for an object resting on an inclined plane in terms of $m$, $g$, and $\theta$.
\part For which $\theta$ is the normal force the biggest? smallest? Does that agree with your answer to (a)?
\part Static friction, like the normal force, provides \emph{as much force as needed} to prevent sliding. For angles before it starts sliding off, what is the magnitude of the normal force, in terms of $m$,$g$,$\theta$?
\part Does the friction force increase or decrease as the plane is tilted higher? Does that agree with your expression in part (c)?
\part Unlike the normal force, static friction has a limit: 
$$
|\vec{f}_s| \leq \mu_s |\vec{N}|.
$$
Given your answers to (b) and (d), explain why there \emph{must} be an angle at which force gravity wins against the friction force and the object slides off of the surface.
\part Show that this angle must be $\theta=\arctan\mu_s$. 
\part My coffee mug slides off my desk if I tilt it at 35$^\circ$. Is $\mu_s$ for the mug-desk surface greater or less than 1? Don't use a calculator.
\end{parts}

\end{questions}




\end{document}
